% User-provided Irodov PDF pp.59-61, original printed pp.50-52.
% Original section 1.7: hydrodynamics, Figs. 1.80-1.86 checked.
\exproblem{315}{51}
理想不可压缩流体沿一条置于水平面内、
横截面积恒定但呈直角弯曲的扁平管稳定流动
（原书图1.80，俯视图）。
图中弯道外侧位置为点1、内侧位置为点2。
问两点压强和速率是否相同，并描述流线疏密。
\exsolution
在弯管内，流体质元须具有指向弯道曲率中心的
法向加速度 $v^2/R_c$。沿弯道横向法线，
稳态欧拉方程近似写作
\[
\frac{\partial p}{\partial n}
=\rho\frac{v^2}{R_c},
\]
其中 $n$ 指向弯道外缘。
因此压强从内缘向外缘升高，
即 $\boxed{p_1>p_2}$。
在图示的扁平弯道中，流线向内弯侧挤密，
给定相邻流管的流量不变，流线越密处局部速度越大，
因而图示 $\boxed{v_1<v_2}$。
流线大体顺着弯道从入口转向出口，
内缘间距较小，外缘较疏。
\excheck 这里的两点位于\emph{不同流线}，
不可不加区分地拿“同一条流线上”的伯努利常数
直接比较 $p_1$ 与 $p_2$。
这一判断与原图中的流线分布一致。

\exproblem{316}{51}
一根水平变截面管道上在截面 $S_1,S_2$
分别安装竖直测压管（原书图1.81），两管水柱高度差
为 $\Delta h$。假设水为理想不可压缩流体、流动稳态。
求管中体积流量 $Q$。
\exsolution
连续性方程给出
$v_1=Q/S_1$、$v_2=Q/S_2$。
两截面等高，沿流线应用伯努利方程：
\[
p_1+\frac12\rho\frac{Q^2}{S_1^2}
=p_2+\frac12\rho\frac{Q^2}{S_2^2}.
\]
测压管显示
$|p_1-p_2|=\rho g\Delta h$。故
\[
g\Delta h=\frac{Q^2}{2}
\left|\frac1{S_1^2}-\frac1{S_2^2}\right|.
\]
解出
\[
\boxed{Q=S_1S_2\sqrt{
\frac{2g\Delta h}{|S_2^2-S_1^2|}}}.
\]
\excheck 两截面必须不等，否则压强差对应流量无法由
该方程唯一求出；若已知
$S_2>S_1$，可直接去掉绝对值为
$\sqrt{2g\Delta h/(S_2^2-S_1^2)}$。

\exproblem{317}{51}
气体在截面积为 $S$ 的管中稳定流动，
中心装有皮托管（原书图1.82）。
差压计液体密度 $\rho_0$，气体密度 $\rho$，
液柱高度差为 $\Delta h$。忽略气体黏性，
求体积流量 $Q$，并说明近似条件。
\exsolution
皮托管迎流孔测停滞压强，
静压孔测流动压强；同高测点的差压为
\[
p_{\rm stag}-p_{\rm static}=\frac12\rho v^2.
\]
U形差压计若两侧气体等密度，其两端压差为
$\Delta p=(\rho_0-\rho)g\Delta h$。
二式相等：
\[
\frac12\rho v^2=(\rho_0-\rho)g\Delta h,
\]
因此
\[
\boxed{Q=Sv=S\sqrt{
\frac{2g\Delta h(\rho_0-\rho)}{\rho}}}.
\]
通常 $\rho_0\gg\rho$，则
\[
\boxed{Q\simeq S\sqrt{\frac{2g\Delta h\rho_0}{\rho}}}.
\]
\excheck 以截面中心速度代替截面平均速度，
还须假定截面内流速分布近似均匀；
实际管流有速度剖面时仅由中心皮托测量不能自动
确定精确总流量。

\exproblem{318}{51}
一个横截面很大的容器中，下层为水、
上层为煤油。
水层厚度 $h_1=30\,\mathrm{cm}$，
煤油层厚度 $h_2=20\,\mathrm{cm}$；
容器底部有小孔通向大气。
忽略黏性，求流出水的速度。
\exsolution
记水和煤油密度分别为 $\rho_1,\rho_2$。
容器很宽，内部两层自由表面下降速度可忽略。
在容器底部小孔处，来自液柱的表压为
\[
p-p_0=\rho_1gh_1+\rho_2gh_2.
\]
在小孔出口，水的静压为大气压 $p_0$。
将液柱压能转化为水的动能：
\[
\frac12\rho_1v^2
=\rho_1gh_1+\rho_2gh_2.
\]
故
\[
\boxed{v=\sqrt{2g\left(
h_1+\frac{\rho_2}{\rho_1}h_2\right)}}.
\]
煤油相对密度取约 $0.80$，得到
$v\simeq\sqrt{2(9.81)(0.30+0.80\cdot0.20)}
=\boxed{3.0\,\mathrm{m/s}}$。
\excheck 不可把上方煤油层直接当成同厚度的水层；
它的贡献须乘密度比 $\rho_2/\rho_1$。

\exproblem{319}{51}
高为 $H=50\,\mathrm{cm}$ 的宽圆柱形水箱装满水，
放在桌面上（箱底即桌面高度）。
在侧壁距底部高度为 $y$ 的位置钻一小孔，
水流水平射出并落到桌面。
忽略黏性，求使落点距箱壁最远的孔高
$y_{\rm opt}$ 与最大水平射程。
\exsolution
自由液面到孔的水头为 $H-y$，
托里拆利公式
$v=\sqrt{2g(H-y)}$。
孔到桌面垂直距离为 $y$，
水平抛体运动的落地时间为
$t=\sqrt{2y/g}$。
射程
\[
L(y)=vt=2\sqrt{y(H-y)},\qquad 0<y<H.
\]
由
$y(H-y)=H^2/4-(y-H/2)^2$
知在 $y=H/2$ 取得全局最大值：
\[
\boxed{y_{\rm opt}=H/2=25\,\mathrm{cm}},\qquad
\boxed{L_{\max}=H=50\,\mathrm{cm}}.
\]
\excheck 孔高接近底部时落地时间太短，
接近水面时喷出速度太小，两端的射程均趋零。

\exproblem{320}{51}
按原书图1.83，弯曲细管迎着水平水流放置，
水流相对细管速率 $v=2.5\,\mathrm{m/s}$。
细管露出水面并在高于水面
$h_0=12\,\mathrm{cm}$ 的闭口端顶部开一个小喷孔。
求喷出的水能够达到喷孔上方的高度 $h$。
\exsolution
迎流管入口将流动动压
$\tfrac12\rho v^2$ 转成停滞压力；
等效水柱高度为
$H_{\rm dyn}=v^2/(2g)$。
管内水从水面上升 $h_0$
消耗重力势能，余下的水头在喷孔处形成喷出速度，
并转为喷孔上方上升高度 $h$。
伯努利方程给出
\[
h+h_0=\frac{v^2}{2g}.
\]
因此
\[
\boxed{h=\frac{v^2}{2g}-h_0
=\frac{(2.5)^2}{2(9.81)}-0.12
\simeq0.20\,\mathrm m=20\,\mathrm{cm}}.
\]
\excheck 必须有 $v^2/(2g)>h_0$
才能从顶端喷出；否则水面仅上升但不出孔。

\exproblem{321}{52}
宽水槽底部开半径为 $R_1$ 的圆孔，
孔上方有半径 $R_2>R_1$ 的封闭圆柱盘，
盘与槽底之间留极窄的水平径向间隙
（原书图1.84）。上部自由水面高于底部 $h$，
液体密度 $\rho$。
求圆柱盘下方间隙内距离圆心 $r$ 处的静压强，
其中 $R_1<r<R_2$。
\exsolution
设外部大气压为 $p_0$。
宽槽静止水体的总水头相当于压强
$p_0+\rho gh$。
在极窄而恒厚度的间隙中，
水沿径向向孔口流动；不可压缩连续性使
$r v_r(r)=$常数。
把小孔边缘 $r=R_1$ 处静压近似取为 $p_0$，
由储水液面的伯努利关系，
其间隙流速大小可取
$v_r(R_1)=\sqrt{2gh}$。
故
\[
v_r(r)=\frac{R_1}{r}\sqrt{2gh}.
\]
对间隙内任意点，
伯努利方程为
$p(r)+\tfrac12\rho v_r(r)^2=p_0+\rho gh$，
于是
\[
\boxed{p(r)=p_0+\rho gh
\left(1-\frac{R_1^2}{r^2}\right)},
\qquad R_1<r<R_2.
\]
\excheck 在 $r=R_1$ 边缘 $p=p_0$，
远离孔口时 $p$ 向静水压 $p_0+\rho gh$ 靠近；
该理想模型忽略极窄缝的黏性损失，
现实中的黏性阻力通常显著。

\exproblem{322}{52}
原书图1.85显示，水平圆筒内有体积 $V$ 的水，
活塞受\emph{恒定}外推力，
水经远小于活塞面积的小孔 $s$ 喷出。
假设无摩擦、无黏性，
要在时间 $t$ 内将水全部挤出，至少需要做多少功？
\exsolution
由水的不可压缩性和题定恒力理想化，
若流量恒定，则
$Q=V/t$，小孔平均喷射速率为
$v=Q/s=V/(st)$。
小孔射出的水动能累积为
\[
K_{\rm out}=
\frac12\rho V v^2
=\frac12\rho V\left(\frac{V}{st}\right)^2.
\]
忽略活塞及筒内水的末态动能、
大气压两侧抵消、没有机械损失，则外力最小功
\[
\boxed{W=\frac{\rho V^3}{2s^2t^2}}.
\]
也可以由所需压差
$\Delta p=\rho v^2/2$ 乘排出体积 $V$ 得同式。
\excheck 恒流量近似要求筒内入口截面远大于小孔截面；
若保留活塞及筒内水的非稳态加速过程，
还需要额外分析瞬态动能与力的时间规律。

\exproblem{323}{52}
一个底面积为 $S$、高度为 $h$ 的竖直圆筒装满水，
底部开面积为 $s\ll S$ 的小孔。
忽略流动黏性，求水全部流出的时间 $\tau$。
\exsolution
设瞬时水面高度为 $z(t)$，
小孔的流出速度由托里拆利关系为
$v(z)=\sqrt{2gz}$。
不可压缩连续性要求箱中体积减少率等于出流量：
\[
-S\frac{\mathrm dz}{\mathrm dt}
=s\sqrt{2gz}.
\]
分离变量积分：
\[
\int_h^0\frac{\mathrm dz}{\sqrt z}
=-\frac{s}{S}\sqrt{2g}\int_0^\tau\mathrm dt.
\]
左侧为 $-2\sqrt h$，
故
\[
\boxed{\tau=\frac{S}{s}\sqrt{\frac{2h}{g}}}.
\]
\excheck $s\ll S$ 是忽略水面下落速度
在伯努利方程中贡献的关键；
实际黏性和孔口收缩会延长排水时间。

\exproblem{324}{52}
长为 $\ell$ 的水平管 $AB$ 绕其一端 $A$
处固定竖直轴以角速度 $\omega$ 匀速旋转
（原书图1.86）。
靠轴端 $A$ 开口，远端 $B$ 封闭但有微小出流孔。
理想液体在管内形成从 $B$ 向 $A$
长度为 $h$ 的液柱，即靠内侧自由界面处于
$x=\ell-h$。求液体相对管道的出流速度。
\exsolution
在随管匀速转动的参考系中，
单位质量的离心势为 $-\omega^2x^2/2$。
液柱内端 $x_A=\ell-h$
与出流孔端 $x_B=\ell$ 都接大气压，
忽略黏性及前端自由界面速度，
沿径向相对流线应用旋转系伯努利积分：
\[
\frac12v^2
=\frac{\omega^2}{2}
\bigl[x_B^2-x_A^2\bigr]
=\frac{\omega^2}{2}\bigl[\ell^2-(\ell-h)^2\bigr].
\]
因此
\[
\boxed{v=\omega\sqrt{2\ell h-h^2}
=\omega h\sqrt{\frac{2\ell}{h}-1}},
\qquad 0<h\le\ell.
\]
\excheck $v$ 受管内液柱长度 $h$ 控制；
题目中的 $h$ 是\emph{水平径向液柱长度}，
不能误解为竖直静水高度。
